\section{Gauss消元与阶梯形}
\label{sec:ls-gaussian-elimination}

后面几节的分解都建立在消元之上。本节先把消元当作手算过程讲清：用第~\ref{sec:ls-elementary-row-operations}节的三种初等行变换把矩阵化为阶梯形，从中读出方程的解，并用同样的过程求逆矩阵。这里只关心精确算术下的结果；怎样把消元过程记录成因子以便复用，以及怎样选主元以控制舍入误差，从第~\ref{sec:ls-lu}节开始讨论。

\subsection{阶梯形与行最简形}
化简时从左到右逐列寻找\term{主元}（pivot）。在当前列尚未使用的行中找一个非零元素；必要时交换行，把它移到当前行，再用它消去同列下方的元素。若该列在余下各行全为零，就跳过这一列，继续向右。

重复到没有可用的行或列为止，得到\term{行阶梯形}（row echelon form）：零行位于底部，每个非零行的首个非零元位于上一行首个非零元的右侧。再把每个主元化为$1$，并消去主元上方的元素，就得到\term{行最简形}（reduced row echelon form）。

阶梯形中主元的个数就是原矩阵的秩。一方面，各步行变换合起来是左乘一个可逆矩阵，不改变矩阵的秩；另一方面，阶梯形的各非零行首元所在的列互不相同、依次右移，这些行线性无关，所以阶梯形的秩等于非零行数，也就是主元个数。第~\ref{sec:ls-solution-cases}节用秩判断方程是否有解，实际计算时就是这样读出秩的。

\subsection{用消元解方程}
解方程$\bm A\bm x=\bm b$时，把右端项作为最后一列附在$\bm A$后面，得到增广矩阵$[\bm A\mid\bm b]$，每次对整行执行同一操作。

三种行变换分别对应交换两个等式、将等式乘以非零数、把一个等式的倍数加到另一个等式上。这些操作不改变解，而且可以撤销，所以新方程组与原方程组的解完全相同。化到阶梯形以后，主元列对应的未知量由方程确定，没有主元的列对应可以任意取值的自由变量。

下面的例子中，第一列的首个元素为零，并且有一列没有主元。
\begin{example}[换行、消元与自由变量]
\label{ex:ls-elementary-row-reduction}
考虑以下增广矩阵。第一列的首个元素为零，先交换前两行：
\[
 \left[\begin{array}{ccc|c}0&1&2&3\\1&2&4&6\\2&4&8&12\end{array}\right]
 \longrightarrow
 \left[\begin{array}{ccc|c}1&2&4&6\\0&1&2&3\\2&4&8&12\end{array}\right].
\]
再把第三行减去第一行的两倍，消去第三行的第一列：
\[
 \left[\begin{array}{ccc|c}1&2&4&6\\0&1&2&3\\2&4&8&12\end{array}\right]
 \longrightarrow
 \left[\begin{array}{ccc|c}1&2&4&6\\0&1&2&3\\0&0&0&0\end{array}\right].
\]
这已是阶梯形。将第一行减去第二行的两倍，得到行最简形
\[
 \left[\begin{array}{ccc|c}1&0&0&0\\0&1&2&3\\0&0&0&0\end{array}\right].
\]
第一、二列是主元列，第三列没有主元。令自由变量$x_3=t$，便读出
\[
 x_1=0,\qquad x_2=3-2t,\qquad x_3=t ;
\]
任意实数$t$都给出原方程组的一个解。

最后的零行
表示$0=0$，不增加约束。若出现系数全为零而右端非零的行，方程组就无解。
\end{example}

这个例子中系数矩阵有两个主元，秩为$2$，比未知量个数少$1$，解集正好带一个自由参数。这与第~\ref{sec:ls-solution-cases}节的结论一致：相容时解集是特解加上零空间，自由参数的个数等于未知量个数减去秩。

\subsection{用消元求逆}
对$n$阶方阵$\bm A$，逆矩阵的第$j$列是方程$\bm A\bm x=\bm e_j$的解（第~\ref{sec:prelim-inverse-matrices}节）。这$n$个方程的系数矩阵相同，消元步骤只由$\bm A$决定，可以一并处理：把$\bm A$与单位矩阵并排，得到$[\bm A\mid\bm I_n]$，对整个增广矩阵作行变换，力求把左侧化为$\bm I_n$。

为什么最后右侧就是逆矩阵？若左侧最终化为$\bm I_n$，各步初等矩阵的乘积记为$\bm E$，它是可逆矩阵，并满足$\bm E\bm A=\bm I_n$；右侧则变为$\bm E\bm I_n=\bm E$。由$\bm E$可逆得$\bm A=\bm E^{-1}$，从而$\bm A\bm E=\bm I_n$。所以右侧的$\bm E$就是$\bm A^{-1}$：
\begin{equation}
 [\bm A\mid\bm I_n]\longrightarrow[\bm I_n\mid\bm A^{-1}].
 \label{eq:ls-gauss-jordan-inverse}
\end{equation}

如果化简时某一列找不到主元，左侧的阶梯形会出现零行，不能化为$\bm I_n$。这时$\bm A$没有逆：若$\bm A$可逆，则$\bm E\bm A$作为两个可逆矩阵之积也可逆；但有零行的矩阵左乘任何矩阵，结果的那一行仍是零，不可能等于$\bm I_n$。

求逆相当于一次解$n$个方程。只需要一个方程的解时，直接对$[\bm A\mid\bm b]$消元即可，先求逆再乘以$\bm b$反而多做运算。
